\documentclass[10pt,a4paper]{article}

\usepackage[T1]{fontenc}
\usepackage[utf8]{inputenc}
\usepackage[margin=1.55cm]{geometry}
\usepackage{microtype}
\usepackage{enumitem}
\usepackage{xcolor}
\usepackage[hidelinks]{hyperref}
\usepackage{titlesec}

\definecolor{ledgerink}{HTML}{19343B}
\setlength{\parindent}{0pt}
\setlength{\parskip}{3.5pt}
\setlist[itemize]{leftmargin=1.25em,itemsep=0pt,topsep=2pt}
\titleformat{\section}{\large\bfseries\color{ledgerink}}{}{0pt}{}
\titlespacing*{\section}{0pt}{5pt}{1pt}
\pagestyle{empty}

\title{\vspace{-1.4cm}\textbf{LedgerCore: Loan Account \& Ledger Service}}
\author{\textit{Author: [Your Name]}}
\date{}

\begin{document}
\maketitle
\vspace{-1.15cm}

\section{Overview and architecture}
LedgerCore is a local loan-servicing backend that creates borrowers and loans, generates
contractual schedules, posts disbursements and repayments, recognizes due interest, and
accrues daily overdue penalties. It is a modular monolith: FastAPI exposes the REST API,
Pydantic validates contracts, SQLAlchemy 2 maps PostgreSQL records, and Alembic manages five
schema revisions. The 13 business tables separate borrowers, snapshotted loan terms,
products, immutable schedules, receipts and allocations, idempotency records, ledger
accounts/journals/entries, accrual identities, demo users, and audit events. Routes own
transactions; accounting, scheduling, security, database, and job code remain separate.

\section{Data model and financial calculations}
All INR amounts are stored as integer paise in \texttt{BIGINT} columns. Python uses
\texttt{Decimal} for intermediate calculations and \texttt{ROUND\_HALF\_UP} to a whole
paisa; binary floating point is rejected for financial inputs. Loans snapshot product terms
so later product edits cannot change an existing schedule. Flat-rate products compute
\(P \times a \times n/m\), while reducing-balance loans use
\(P r(1+r)^n/((1+r)^n-1)\), with a safe zero-rate branch. Deterministic remainder handling
keeps scheduled principal exactly equal to original principal. A schedule is a contractual
plan; principal, recognized interest, and penalties become receivables only through posted
ledger journals. Consequently, the current balance is not automatically a future payoff quote.

\section{Double-entry accounting}
Each posted journal has at least two immutable entries and equal debit and credit totals.
Disbursement debits loan principal receivable and credits cash. Scheduled interest recognition
debits interest receivable and credits interest income. Repayment debits cash and credits the
recognized penalty and interest receivables, then due or future principal according to policy.
The system stores no mutable authoritative balance: read models sum posted entries, including
supported repayment reversals. PostgreSQL constraints and triggers supplement service checks
by rejecting unbalanced posting and mutation of posted financial history.

\section{Idempotency and concurrency}
Repayment requests require actor/operation-scoped idempotency keys. A canonical payload
fingerprint distinguishes an exact retry, which returns the original successful response,
from a changed request, which returns a conflict. Source plus external-event identity provides
independent duplicate detection. Each financial operation uses one caller-owned database
transaction for recognition, receipt, allocations, journal entries, idempotency result, and
audit. \texttt{SELECT FOR UPDATE} serializes mutations on a loan, while uniqueness races are
retried after rollback. If a response is lost after commit, retrying the frozen request recovers
the original receipt rather than posting again.

\section{Replay and reconciliation}
The retained deterministic replay (seed 20261114) contains 250 invented borrowers, 400 loans,
three products, and exactly 5,000 repayment delivery attempts through the HTTP API. Observed
results were 4,510 new repayments, 330 duplicate responses, 160 rejected deliveries, and 250
accepted out-of-order events. Reconciliation counted 9,903 posted journals and
3,641,480,368 paise on each side of the ledger, with zero unbalanced journals and no reported
principal, interest, or penalty differences. The scenario manifest was generated independently
of actual API outcomes; these figures establish consistency for this dataset rather than proof
of every possible financial edge case.

\section{Limitations and next step}
\texttt{X-Demo-User} is development-only impersonation, not verified production identity;
there is no tenant, borrower-ownership, or officer-assignment isolation. Reversal support is
limited to repayment journals, and Docker/Compose execution was not verified in the final
environment (native PostgreSQL was). The most valuable next step is verified OIDC identity
with explicit tenant-scoped authorization and least-privilege database roles, followed by
deployment-level operational testing.

\end{document}
